\section{部署蓝图：技术栈与组织协同}

\subsection{参考技术栈}
结合课程观点，一个可落地的 Agent 栈可以分为六层：接口层、编排层、记忆层、工具层、验证层、观测层。每层都应具备最小自治能力，并通过稳定协议通信。

\begin{table}[H]
\centering
\begin{tabular}{p{2.4cm}p{4.0cm}p{4.6cm}}
\toprule
\textbf{层级} & \textbf{核心职责} & \textbf{关键产物} \\
\midrule
接口层 & 接收任务、权限校验、输出格式化 & Task spec、权限上下文 \\
编排层 & 任务分解、路由、重规划触发 & Action graph、执行计划 \\
记忆层 & 分层读写、压缩、淘汰、冲突处理 & 记忆索引与摘要 \\
工具层 & 统一工具协议、错误封装、重试 & 可调用 API catalog \\
验证层 & 终局与步骤验证、风险拦截 & 验证报告与风险标签 \\
观测层 & 轨迹追踪、指标聚合、告警 & Trace、Dashboard、审计日志 \\
\bottomrule
\end{tabular}
\caption{Agent 系统部署的六层参考结构}
\end{table}

\subsection{组织协同模型}
Agent 项目常见组织问题是：模型组、平台组、业务组各自优化局部目标。课程启发是建立 ``任务成功率'' 为统一北极星指标，并让不同团队共享同一套轨迹数据和失败分类。

\begin{importantbox}{组织协同的三条实践}
\begin{enumerate}
    \item 统一事件 schema，避免各组口径不一致；
    \item 统一回滚机制，确保线上风险可控；
    \item 统一优先级框架，以高损失失败优先修复。
\end{enumerate}
\end{importantbox}

\subsection{治理与合规}
在企业部署中，Agent 不只是技术问题，还涉及权限最小化、敏感数据处理、审计留痕和责任归属。课程建议在系统层内置治理能力，而不是上线后补丁式合规。

\begin{warningbox}{治理能力必须前置}
\begin{itemize}
    \item 没有权限分层，工具调用会出现越权风险；
    \item 没有审计日志，事故后难以完成责任追踪；
    \item 没有红线规则，模型异常可能快速扩散。
\end{itemize}
\end{warningbox}

\begin{knowledgebox}{部署前检查清单}
\begin{itemize}
    \item 是否具备任务级回滚与人工接管？
    \item 是否具备步骤级日志与可复放轨迹？
    \item 是否具备预算闸门与故障熔断？
    \item 是否完成敏感场景红队测试？
\end{itemize}
\end{knowledgebox}

\subsection{本章小结}
Agent 部署是 ``技术 + 组织 + 治理'' 的联动工程。只有将三者合并设计，系统才能稳定穿越从 Demo 到生产的断层。

